\section{Introduction}

\begin{sloppypar}
Web agents, or LLM-based systems that complete tasks by directly browsing and interacting with live websites, are increasingly deployed in retrieval-augmented assistants, deep-research tools, and enterprise automation pipelines. Agents such as WebVoyager~\cite{webvoyager}, AgentOccam~\cite{agentoccam}, and BrowserUse~\cite{browseruse} now underpin a fast-growing class of products that resolve queries no static corpus can answer: navigating to authenticated pages, completing multi-step lookups across catalogs, and extracting structured information from sites the user does not need to visit themselves. This capability far exceeds what stateless retrieval pipelines (e.g., Perplexity, ChatGPT browsing) achieve by querying search indices and synthesizing from snippets -- they cannot maintain authenticated state, complete multi-step trajectories, or reach content beyond search coverage -- but it comes at substantial overhead: per-task latencies run 30-120 secs and API costs run \$0.20-0.50, 1-2 orders of magnitude more than stateless retrieval. %At deployment scale, these per-task costs compound into operational expenses that bound adoption.
\end{sloppypar}

The high overheads of web agents trace to three off-the-shelf components, each designed for problems broader than the web tasks they typically execute. Frontier LLMs are trained for long-context multi-hop reasoning; browsers are engineered for arbitrary and secure human interaction with any web page; ReAct is a general decision-making framework for agents in unconstrained environments with sparse feedback. Each component is highly capable, but fails to exploit the structural regularities that the web exhibits. The first such property is \emph{step heterogeneity}. Most steps in a typical web agent trajectory are mechanical (fetch a known URL, locate a labeled field on a page, click a stable navigation element) and do not require frontier reasoning or full browser execution. The second is \emph{task-driven structure}. Sites are purpose-built around the queries users bring to them, so URL templates, page layouts, and answer formats are stable across queries of the same type. For instance, the navigation pattern that resolves one arXiv lookup resolves the next, with only query-specific identifiers changing. Together, these properties imply that many steps can run on cheaper components, and many can be skipped entirely by following trajectories the site has already revealed.

Our analysis of popular web agent tasks confirms that the opportunity is substantial (\S\ref{sec:opp}). Hand-engineered programs that exploit website structure by using trajectory knowledge to fetch via direct URLs, prune to answer-bearing content, and extract with cheaper models can run 66.7-94.9\% faster and 17.7-100.7$\times$ cheaper than the off-the-shelf ReAct agent without any accuracy loss. However, automating this strategy in practice is difficult; naively substituting these cheaper techniques inside an off-the-shelf ReAct loop drops average task success by 60\%. The issue is that cheap models struggle to choose the right action from the hundreds of interactive elements on a typical page, and direct-URL fetching requires knowing in advance precisely which page will contain the answer for a given query (\S\ref{sec:chal}). Worse, errors compound across a trajectory, so a single wrong step can invalidate everything downstream. The core question is thus how to identify, when a query arrives, which shortcuts preserve accuracy for \emph{that} query -- a question made difficult by the heterogeneity of tasks, the variation across pages in a site, and the per-step diversity in resource needs.

%Our analysis of popular web agent tasks confirms the opportunity is substantial (\S\ref{sec:motivation:handtune}). Hand-engineered programs that exploit website structure by using trajectory knowledge to fetch via direct URLs, prune to answer-bearing content, and extract with cheaper models can drop per-task cost by X$\times$ and latency by Y\% without accuracy loss. However, realizing this strategy in practice is difficult --  naively substituting these cheaper techniques inside an off-the-shelf ReAct loop drops task success by Y\%, because even a single wrong step can compound across the trajectory and invalidate the result (\S\ref{sec:motivation}) \mw{Do we want to mention the other core reasons for failure? Namely that the action space is too large for small models, direct fetching requires a priori knowledge of trajectories, etc. }. The core question is thus how to identify, when a query arrives, which shortcuts preserve accuracy for \emph{this} query? Tasks vary in objects and objectives, sites vary internally across page types, and the minimum-sufficient resource varies per step within a single trajectory, all of which make this identification difficult to do cheaply and reliably.

To address these challenges, we present \textbf{\sys}, a drop-in acceleration framework for existing web agents. The key insight underlying \sys{} is that web sites are \emph{predictable enough to precompute} navigation paths offline, and \emph{structured enough to check} answer correctness online. Sites change rarely relative to query rates, so trajectory patterns characterized once per task type remain valid across many future queries. Further, since sites impose stable formats on task result pages (e.g., search results), a fast path's output can be easily validated against both the query and the expected schema, without invoking the expensive agent. Building on this, \sys{} profiles each reachable site offline, recording for each task type both the shortcuts that apply (URL patterns reaching its result pages) and what correct outputs look like (the answer schemas of those pages). At runtime, when a query arrives, \sys{} consults the profile and, if covered, speculatively attempts a fast path via direct URL fetch and cheap model processing, accepting the result only when a small judge model deems it consistent with the query and the recorded schema. Upon failure, \sys{} falls back to the full ReAct agent, warm-started at the URL the fast path reached -- since trajectories share predictable prefixes, even failed shortcuts reduce the work that the agent must redo.




%To address these challenges, we present \textbf{\sys}, a drop-in acceleration framework for existing web agents. The key insight underlying \sys is that web sites are cheap to \emph{profile} offline and cheap to \emph{verify} online. Sites change rarely relative to query rates, so an infrequent investment in modeling each site's recurring trajectories amortizes across many future queries. Moreover, since answer formats are determined by site purpose, a correct output can be recognized without re-running the full agent. Building on this, \sys{} profiles each reachable site offline, recording URL templates and answer schemas for the task types it encounters. At runtime, when a query arrives, \sys{} consults the profile and, if covered, attempts a fast path that fetches the candidate URL, extracts with a non-frontier model, and validates the output against the recorded schema. Upon failures, \sys{} falls back to the full ReAct agent, warm-started at the URL the fast path already reached, so navigational progress carries over rather than being discarded.


We evaluate \sys on two representative benchmarks (WebVoyager, WebShop) across three popular ReAct-based web agents (BrowserUse, AgentOccam, WebVoyager). We consider two deployment modes that exploit the same underlying mechanism. In \emph{accelerate} mode, \sys lowers median per-task cost by 1.9$\times$ and latency by 33.4\% while preserving accuracy. In \emph{aggregate} mode, \sys reinvests the saved compute budget into multiple trials that explore different trajectories with verifier-selected outputs. Within the single-trial baseline's cost budget, this lifts end-to-end accuracy by up to 16.7 percentage points (4.2 pp with majority vote).

